%!TEX root = ../thesis.tex

\chapter{Conclusion}\label{chap:conclusion}

We have implemented a fuzzing framework based on the EfficientZero deep reinforcement
learning algorithm and the MOGI emulator.
The targets we evaluated show that it is possible to learn a useful input sequence
in the input-action model as well as useful mutation operators in the corpus-action model.
Coverage climbs for most targets and produces recognizable structure for the cJSON
and libxml2 targets.
We also analyzed drawbacks that are mainly concerned with throughput.

We answer the research questions posed in~\cref{sec:introduction-rqs} as follows:

\begin{enumerate}[label=\textbf{RQ\arabic*.}, ref=RQ\arabic*, leftmargin=*]
  \item\label{rq:learnability}
    \enquote{Can EfficientZero learn a useful policy on fuzzing-shaped rewards?
    Coverage rewards are sparse, and observations alias distinct program states
    while giving only a partial view of the true state.}

    We answer this question with a clear yes.
    The targets can be successfully learned by using a fuzzing-shaped reward
    derived from coverage instrumentation.

  \item\label{rq:action-space}
    \enquote{How should the fuzzing action space be modeled?
    We pursue two approaches: modeling actual inputs a program expects as actions,
  as well as modeling mutation operators acting on a corpus as actions.}

  \item\label{rq:planning}
    \enquote{Does planning with a learned model help? A core part of EfficientZero is
    \enquote{imagining} how a program behaves given an input. How well can this be learned?}
\end{enumerate}

Whether this deep RL-based approach can ultimately beat conventional \gls{cgf} fuzzers
is not clear.
Trying to act \enquote{smarter} sounds good in theory and also works for toy examples,
but the major advantage of \gls{cgf}, namely the high execution throughput to test
millions of seeds in a short time, is hard to beat.
In addition, our approach suffers from the same major drawback of not being able
to pass checksum and magic byte checks, as well as RL-specific problems like policy collapse.
The potential in our approach lies in learning regularity in PUT behavior and exploiting
this in the fuzzing process, as well as being able to flexibly plug in new observation
and action models and steer exploration using reward shaping.
Considering the open points in~\cref{chap:future-work}, we think various
improvements are still open.
Better engineering of a solution might decrease the execution throughput gap between ML-based
solutions and traditional \gls{cgf}, while incorporating new research might increase
an agent's ability to exploit available information.
